\documentclass[conference]{IEEEtran}
\IEEEoverridecommandlockouts

\usepackage{cite}
\usepackage{amsmath,amssymb,amsfonts}
\usepackage{graphicx}
\usepackage{textcomp}
\usepackage{xcolor}
\usepackage{booktabs}
\usepackage{array}
\usepackage{url}

\newcolumntype{L}[1]{>{\raggedright\arraybackslash}p{#1}}
\newtheorem{theorem}{Theorem}
\newtheorem{proposition}{Proposition}

\def\BibTeX{{\rm B\kern-.05em{\sc i\kern-.025em b}\kern-.08em
    T\kern-.1667em\lower.7ex\hbox{E}\kern-.125emX}}

\begin{document}

\title{Security Limits of Adaptive Uncertainty in Edge Digital Twins for Mobile Rovers}

\author{\IEEEauthorblockN{Shrey [Last Name]}
\IEEEauthorblockA{\textit{[Department]} \\
\textit{[University]}\\
[City], USA \\
[email]}
}

\maketitle

\begin{quote}
\textbf{Artifact status (July 29, 2026):} The numerical results in this draft
were generated before the independent security-predictor revision. They are
retained as reproducible historical results, but the matched alarm thresholds,
attack summaries, tables, and figures must be regenerated before submission.
The existing rover logs are sufficient for that software-only rerun.
\end{quote}

\begin{abstract}
Edge-hosted digital twins can monitor low-cost mobile robots with richer state
estimation and security logic than the robot can execute locally. However, the
same uncertainty estimates that make the detector usable under wheel slip,
terrain variation, GPS noise, and network delay can also change attack
detectability. This creates a security question that is usually hidden by
fixed-covariance evaluations: can adaptive uncertainty itself become an attack
surface? This paper studies semantic GPS-coordinate manipulation against
a Waveshare UGV01 tracked rover with BN220 GPS, encoder, IMU, and timestamped
telemetry. Using real rover logs, we evaluate replay-injected GPS attacks
against fixed, robust, sequential, and adaptive detector variants at matched
run-level false-alarm calibration. The frozen campaign contains 20 physical
runs, 24 attack profiles, three attack start times, 1,440 unique
attack-run-start combinations, and 15,840 detector-run evaluations; an expanded
replay grid increases this to 44 attack profiles and 29,040 detector-run
evaluations. Results show that large step attacks are detectable, but slow GPS
drift remains difficult under the current platform uncertainty. For the
representative 0.05 m/s cross-track drift, all variants have zero run-level
detection probability at the locked operating point, while naive
residual-coupled adaptation has 0.150 tolerance-exceeding paired divergence
before alarm at the 5 m tolerance. Fixed and robust baselines are 0.017 under
the same metric. The paper contributes a threat model, detectability analysis,
baseline suite, post-campaign paired comparisons, and real-log evidence package
for studying the security limits of adaptive uncertainty in edge digital
twins. The present results are an offline real-log study; independent physical
ground truth and prospective live validation remain future work.
\end{abstract}

\begin{IEEEkeywords}
Digital twin, mobile robot security, GPS spoofing, semantic telemetry attacks,
extended Kalman filter, adaptive covariance, edge computing.
\end{IEEEkeywords}

\section{Introduction}

Low-cost mobile robots increasingly offload estimation, monitoring, and
decision support to a nearby edge workstation \cite{ref48,ref61}. This architecture is attractive
because a small embedded controller may not have the compute budget for
multiple estimators, attack injection, statistical replay, and full telemetry
logging. Yet the edge-hosted design also creates a measurement-security
problem: the detector must distinguish malicious GPS-coordinate manipulation
from ordinary uncertainty caused by GPS noise, tracked-drive slip, terrain
disturbance, timing jitter, and kinematic mismatch \cite{ref11,ref43,ref65}.

This paper studies that problem on a physical UGV01 rover. The rover streams
combined telemetry including GPS, encoder counts, IMU values, voltage,
sequence metadata, and source and edge timing fields. The edge digital twin
replays the real telemetry through an extended Kalman filter (EKF), injects
controlled semantic GPS attacks, and compares multiple detector families under
matched benign false-alarm calibration.

The central question is not simply whether a detector can identify an obvious
GPS jump. Instead, we ask when an edge-hosted digital twin has enough
statistical sensitivity to detect mission-relevant semantic coordinate
manipulation, and whether an adaptive uncertainty mechanism can itself become
an attack surface. In particular, if the covariance predictor learns from a
residual containing attacker-controlled GPS data, a slow drift can be
misinterpreted as increased benign uncertainty. This increases the innovation
covariance and reduces normalized innovation evidence.

The contributions of this work are:
\begin{itemize}
    \item a real UGV01 telemetry and replay pipeline for edge-hosted
    digital-twin security evaluation;
    \item a coordinate-only semantic GPS threat model with protected encoder,
    IMU, sequence, command, and edge-time evidence;
    \item a matched false-alarm comparison of 11 detector and estimator
    variants, including fixed NIS, raw residuals, robust innovation gating,
    Huber EKF, CUSUM, and adaptive covariance variants;
    \item an offline attack campaign covering 1,440 frozen unique
    attack-run-start combinations, 15,840 frozen detector-run evaluations,
    and a 29,040-evaluation expanded replay grid;
    \item empirical evidence, paired detector differences, gate-behavior
    summaries, and runtime estimates showing that slow semantic drift is
    difficult under the current platform uncertainty and that naive adaptive
    covariance can increase tolerance-exceeding paired divergence before
    alarm.
\end{itemize}

\section{Background and Related Work}
\label{sec:related_work}

This work sits at the intersection of mobile-robot localization, localization
security, and edge-hosted cyber-physical monitoring. EKF-based localization is
standard for low-cost mobile robots because it can combine a kinematic motion
model with noisy exteroceptive measurements. In tracked or differential-drive
rovers, however, the motion model is only approximate: slip, surface changes,
turning bias, and encoder quantization alter the relation between commanded
motion and physical trajectory \cite{ref11}. A detector that ignores this
uncertainty may false alarm during ordinary motion; a detector that overstates
uncertainty may miss attacks.

Localization security has studied spoofing, false-data injection, and secure
state estimation in cyber-physical systems \cite{mo2010,pasqualetti2013,fawzi2014}.
Many of these results assume a fixed model or a stationary covariance. That is
appropriate for some process-control settings, but it is incomplete for a
small rover whose uncertainty changes with surface, speed, GPS quality, and
packet timing. IoT localization surveys and sensor-spoofing defenses further
show that syntactically valid position messages can still be semantically
malicious \cite{ref43,ref60}. This paper focuses on that semantic layer: the
packet may be well formed, but the GPS coordinate can be wrong.

Edge computing changes the security problem again. Offloading makes richer
replay, logging, and detection feasible, but edge arrival time, queueing, and
source-to-estimate latency become part of the operating condition
\cite{ref48,ref61}. A mobile edge digital twin therefore needs both
state-estimation uncertainty and timing uncertainty. The novelty of this work
is not a new EKF by itself, nor a generic GPS-spoofing detector. The specific
question is whether adaptive uncertainty estimation, used to remain calibrated
under changing physical and network conditions, can become a feedback channel
that an attacker exploits.

\begin{table*}[t]
\centering
\caption{Positioning relative to adjacent research areas. A check mark denotes a central evaluation focus rather than incidental mention.}
\label{tab:novelty}
\scriptsize
\begin{tabular}{L{0.23\linewidth} c c c c c c}
\toprule
\textbf{Area} & \textbf{Real rover logs} & \textbf{Slow drift} & \textbf{Adaptive covariance security} & \textbf{Matched FA} & \textbf{Edge timing} & \textbf{Ground truth} \\
\midrule
Classical secure state estimation & -- & partial & -- & partial & -- & -- \\
GPS/GNSS spoofing detection & partial & partial & -- & partial & -- & partial \\
Robust residual detection & partial & partial & -- & partial & -- & partial \\
Digital-twin security studies & partial & partial & -- & -- & partial & -- \\
This offline real-log study & yes & yes & yes & yes & yes & no \\
Target prospective extension & yes & yes & yes & yes & yes & yes \\
\bottomrule
\end{tabular}
\end{table*}

\section{Research Questions}
\label{sec:research_questions}

The study is organized around five questions.

\textbf{RQ1: Detectability.} What GPS step magnitude is required to obtain a
target run-level detection probability under matched false-alarm calibration?

\textbf{RQ2: Slow drift.} Are gradual coordinate manipulations harder to
detect than abrupt offsets on the same real rover logs?

\textbf{RQ3: Covariance poisoning.} Does residual-coupled adaptive covariance
inflate $Q_k$ or $S_k$ under attack and suppress normalized innovation
evidence?

\textbf{RQ4: Baseline comparison.} Do ordinary robust or sequential baselines
remove the observed vulnerability at the same false-alarm operating point?

\textbf{RQ5: Defense boundary.} Do GPS-independent and evidence-gated
adaptation reduce the direct attacked-GPS-to-covariance feedback path, and
what claims remain unsupported without independent ground truth?

\section{Threat Model}
\label{sec:threat_model}

The defended system is an edge-hosted digital twin monitoring a tracked UGV01
rover. The rover reports GPS coordinates, encoder counts, IMU readings,
voltage, command state, sequence numbers, and timing metadata. The edge records
arrival, queue, estimate, and alarm timestamps using its own monotonic clock.

The primary adversary is a white-box, coordinate-only semantic GPS attacker.
The attacker may alter the GPS coordinate values supplied to the estimator and
may choose the attack direction, magnitude, rate, and start time. The attacker
knows the estimator, detector, thresholds, and covariance policy. However, the
attacker cannot modify motor commands, encoder counts, IMU values, GPS-quality
fields, sequence values, edge-observed arrival times, detector implementation,
alarm outputs, or the stored raw logs.

This threat model represents the detector-visible consequence of GPS
coordinate spoofing or semantic telemetry manipulation \cite{ref2,ref60}. It does not claim an
over-the-air RF spoofing implementation. It also does not model full edge-host
compromise; an adversary that can alter the detector code, thresholds, alarm
path, or logs could trivially defeat this experimental defense.

\begin{table}[t]
\centering
\caption{Attack families and evaluation role.}
\label{tab:attack_families}
\scriptsize
\begin{tabular}{L{0.13\linewidth} L{0.36\linewidth} L{0.37\linewidth}}
\toprule
\textbf{Family} & \textbf{Attacker controls} & \textbf{Protected evidence} \\
\midrule
A1 semantic GPS content & GPS coordinates and injection schedule &
Encoder, IMU, command, sequence, GPS quality, edge time, detector, alarm \\
A2 delivery manipulation & Packet delay, loss, order, or repetition &
Payload values and detector implementation under replay assumption \\
A3 combined manipulation & GPS content and delivery schedule &
Protected estimator, alarm, and evaluation reference \\
\bottomrule
\end{tabular}
\end{table}

The main results in this paper evaluate A1. Controlled delay and jitter are
used as edge-side operating conditions and robustness checks, not as evidence
of malicious intent.

\section{System and Telemetry}
\label{sec:system}

\subsection{Rover and Firmware}

The physical platform is a Waveshare UGV01 tracked rover with an ESP32
controller, wheel encoders, onboard IMU, voltage sensing, and a BN220 GPS
receiver. Low-level firmware path names, wiring details, and operator commands
are part of the reproducibility package rather than the central scientific
claim.

The firmware exposes two relevant telemetry commands. Command \texttt{T:146}
returns GPS-only telemetry. Command \texttt{T:147} returns combined telemetry,
including sequence number, source timestamp, encoder counts, IMU fields,
voltage, GPS validity, latitude, longitude, satellite count, HDOP, GPS speed,
and firmware-side checksum counters. The edge logger wraps this packet with
send time, edge arrival time, queue depth, estimate time, alarm time, HTTP
latency, packet-loss indicators, and stale-packet indicators.

\subsection{Edge Digital Twin}

The edge pipeline parses telemetry, converts GPS latitude and longitude into a
local frame, predicts rover motion from encoder-derived controls, and scores
GPS measurements against an EKF innovation. The security branch is propagated
from protected non-GPS evidence and edge timing. The attacked GPS measurement
is evaluated against the pre-update prediction. During attack evaluation, GPS
coordinates are never used as the only source of the security reference.

\begin{figure}[t]
\centering
\fbox{\begin{minipage}{0.92\linewidth}
\centering
\scriptsize
UGV01 sensors and firmware\\
GPS, encoders, IMU, voltage, sequence, source time\\[2pt]
$\downarrow$ Wi-Fi telemetry\\[2pt]
Edge logger\\
arrival time, queue depth, packet health, clock offset\\[2pt]
$\downarrow$\\[2pt]
\begin{tabular}{c c}
Clean replay branch & Attacked GPS branch \\
protected controls/IMU/encoders & semantic GPS injection
\end{tabular}\\[2pt]
$\downarrow$\\[2pt]
EKF prediction, covariance policy, detector score, persistent alarm
\end{minipage}}
\caption{Edge-hosted digital-twin evaluation architecture. The primary A1 attacker changes GPS coordinates supplied to replay, while encoder, IMU, sequence, and edge-time evidence remain protected.}
\label{fig:architecture}
\end{figure}

\section{Detector and Uncertainty Model}
\label{sec:detector}

Let the rover state be
\begin{equation}
    x_k = [p_{x,k}, p_{y,k}, \theta_k]^T,
\end{equation}
with nonlinear transition
\begin{equation}
    x_{k+1} = f(x_k, u_k, \Delta t_k) + w_k,
    \qquad w_k \sim \mathcal{N}(0,Q_k),
\end{equation}
and GPS measurement
\begin{equation}
    z_k = Hx_k + v_k + a_k,
    \qquad v_k \sim \mathcal{N}(0,R_k).
\end{equation}
Here $a_k$ is the semantic GPS attack vector. The innovation and innovation
covariance are
\begin{equation}
    \nu_k = z_k - H\hat{x}_k^-,
    \qquad
    S_k = HP_k^-H^T + R_k.
\end{equation}
The primary detector score is the normalized innovation squared (NIS):
\begin{equation}
    d_k = \nu_k^T S_k^{-1}\nu_k.
\end{equation}
The operational alarm is motion-gated and persistent: the main NIS-family
detectors alarm when three of the latest five monitored scores exceed the
benign-locked threshold. GPS-jump and CUSUM baselines use their own
variant-specific score and threshold.

\subsection{Oracle Detectability Bound}

For a single GPS innovation under a fixed covariance, the attack-detection
problem can be written as a Gaussian mean-shift test:
\begin{equation}
    H_0:\nu_k\sim\mathcal{N}(0,S_k),
    \qquad
    H_1:\nu_k\sim\mathcal{N}(a_k,S_k).
\end{equation}
When the attack direction is known, the Neyman--Pearson optimal statistic is
\begin{equation}
    T_k = \frac{a_k^T S_k^{-1}\nu_k}{\sqrt{\lambda_k}},
    \qquad
    \lambda_k = a_k^T S_k^{-1}a_k.
\end{equation}
Then $T_k|H_0\sim\mathcal{N}(0,1)$ and
$T_k|H_1\sim\mathcal{N}(\sqrt{\lambda_k},1)$. For per-update false-alarm
probability $\alpha$, the best achievable single-update detection probability
for the known-direction attack is
\begin{equation}
    P_D^{\mathrm{oracle}} =
    1-\Phi\left(\Phi^{-1}(1-\alpha)-\sqrt{\lambda_k}\right).
\end{equation}
Thus a target detection probability $\beta$ requires
\begin{equation}
    \lambda_k \geq
    \left[\Phi^{-1}(1-\alpha)+\Phi^{-1}(\beta)\right]^2.
\end{equation}
This bound is not the implemented detector; it is a reference for
interpreting poor detection. If the empirical detector misses small attacks
whose oracle $\lambda_k$ is also small, the limitation is partly statistical.
If the oracle predicts detectability but the implemented detector fails, the
limitation lies in the detector, persistence rule, or covariance model.

\subsection{Three-of-Five Run-Level Alarm}

The implemented alarm is not a one-update decision. If a monitored update has
threshold-crossing probability $p$ and if crossings are momentarily treated as
independent, one five-sample window alarms with probability
\begin{equation}
    P_{3/5}(p)=
    \sum_{j=3}^{5}\binom{5}{j}p^j(1-p)^{5-j}.
\end{equation}
The actual run-level probability is estimated empirically because windows
overlap and innovations are serially correlated. This approximation is still
useful: it explains why the operational detector can be conservative even
when a single update occasionally exceeds threshold. It also motivates
reporting detection delay and tolerance-exceeding paired divergence before
alarm in addition to per-update NIS.

\subsection{Covariance-Poisoning Mechanism}

If covariance adaptation depends on a residual containing attacker-controlled
GPS data, then a slow drift can increase the proposed process covariance. The
following local result states why that matters for detectability.

\begin{proposition}[Innovation covariance and normalized attack energy]
\label{prop:covariance_order}
Let $S_1$ and $S_2$ be positive definite innovation covariances with
$S_2 \succeq S_1$. For any fixed attack vector $a$,
\begin{equation}
    a^T S_2^{-1} a \leq a^T S_1^{-1} a.
\end{equation}
\end{proposition}

Thus, once covariance inflation reaches the innovation covariance, the same
physical attack has lower normalized statistical energy. This does not by
itself prove every nonlinear EKF trajectory is globally ordered, but it
identifies the mechanism measured in the paired replay study.

\begin{theorem}[Scoped covariance monotonicity for paired replay]
\label{thm:dynamic_poisoning}
Consider a linear Kalman filter covariance recursion, or a locally linearized
EKF replay in which the paired clean and attacked branches use identical
initial covariance, fixed linearization matrices $F_k$ and $H_k$, and identical
measurement covariance $R_k$. If the attacked adaptive process covariance
satisfies
\begin{equation}
    Q_k^a = Q_k^0 + \Delta Q_k, \qquad \Delta Q_k \succeq 0,
\end{equation}
for all $k$, then the predicted covariance and innovation covariance satisfy
\begin{equation}
    P_k^{a,-}\succeq P_k^{0,-},
    \qquad
    S_k^a\succeq S_k^0.
\end{equation}
Consequently, for a fixed attack vector $a_k$, the normalized attack energy
satisfies
\begin{equation}
    a_k^T(S_k^a)^{-1}a_k
    \leq
    a_k^T(S_k^0)^{-1}a_k.
\end{equation}
\end{theorem}

\emph{Proof sketch.} The base case follows from identical initial covariance.
If $P_k^a\succeq P_k^0$, then prediction gives
$F_kP_k^aF_k^T+Q_k^a\succeq F_kP_k^0F_k^T+Q_k^0$. The Kalman covariance update
with fixed $H_k$ and $R_k\succ0$ is monotone in the prior covariance, so the
posterior ordering is preserved by induction. Since
$S_k=H_kP_k^-H_k^T+R_k$, innovation covariance ordering follows, and
Proposition~\ref{prop:covariance_order} gives the normalized-energy result.
This theorem is a scoped covariance result for linear filters and fixed
linearization paired replay. It is not claimed as a global theorem for the
complete nonlinear EKF trajectory.

The theorem gives the causal chain:
\[
    \Delta Q \rightarrow \Delta P^- \rightarrow \Delta S
    \rightarrow \text{lower normalized attack energy}.
\]
This is the theoretical basis for the covariance-poisoning experiment.

\subsection{Evidence-Gated Adaptation}

The evidence-gated variant allows GPS-residual-driven adaptation only when
non-GPS evidence independently supports increased uncertainty. The protected
evidence sources in the current implementation are vertical IMU acceleration,
yaw-rate magnitude, edge-observed timing mismatch, and stale-packet status.
GPS coordinates are not used to open the gate. Define
\begin{equation}
\begin{split}
e_k^{\mathrm{imu}} &=
\mathbb{1}\{|a_{z,k}-g|\geq 0.8\ \mathrm{m/s^2}\}
\vee
\mathbb{1}\{|\omega_{z,k}|\geq0.35\ \mathrm{rad/s}\},\\
e_k^{\mathrm{time}} &=
\mathbb{1}\{|\Delta t_k^{\mathrm{edge}}-\Delta t_k^{\mathrm{src}}|
\geq0.20\ \mathrm{s}\},\\
e_k^{\mathrm{stale}} &= \mathbb{1}\{\text{packet is stale}\}.
\end{split}
\end{equation}
The residual-feedback gate is
\begin{equation}
    b_k =
    \mathbb{1}\{d_{k-1}\leq \gamma\}
    \wedge (e_k^{\mathrm{imu}}\vee e_k^{\mathrm{time}})
    \wedge \neg e_k^{\mathrm{stale}},
\end{equation}
where $\gamma$ is the locked detector threshold for the variant. The first term
prevents residual feedback after an already suspicious innovation; the second
requires independent IMU or timing evidence; the third rejects stale packets.

When $b_k=1$, the covariance proposal is accepted; when $b_k=0$, adaptation is
frozen:
\begin{equation}
Q_k =
\begin{cases}
\phi_Q Q_{k-1} + (1-\phi_Q)\widetilde{Q}_k, & b_k=1,\\
Q_{k-1}, & b_k=0.
\end{cases}
\end{equation}
If $Q_{\min}\preceq \widetilde{Q}_k \preceq Q_{\max}$ and $N_{\mathrm{pass}}$
updates pass the gate, then
\begin{equation}
Q_k \preceq
\phi_Q^{N_{\mathrm{pass}}}Q_0+
(1-\phi_Q^{N_{\mathrm{pass}}})Q_{\max}.
\end{equation}
The gate therefore does not guarantee attack detection, but it bounds the
additional covariance cover that can accumulate through accepted adaptive
updates. This is the intended security role of evidence gating.

\subsection{Attack-to-Adaptation Feedback Loop}

The residual-coupled adaptive variant creates the following influence path:
\begin{equation}
    a_k \rightarrow \nu_k \rightarrow \widehat{Q}_k
    \rightarrow P_k^- \rightarrow S_k \rightarrow d_k.
\end{equation}
This path is absent in fixed covariance, reduced in GPS-independent
adaptation, and conditionally limited in evidence-gated adaptation. If the
adaptive proposal satisfies a bounded-influence condition
\begin{equation}
    0 \preceq \widehat{Q}_k-Q_k^0 \preceq c_Q\|a_k\|^2 I,
\end{equation}
and at most $N_{\mathrm{pass}}$ attack-influenced updates pass the gate, then
the accumulated attack-induced process covariance is bounded by
\begin{equation}
    Q_k-Q_k^0 \preceq
    (1-\phi_Q^{N_{\mathrm{pass}}})c_Q A_{\max}^2 I,
\end{equation}
where $A_{\max}$ bounds the attack magnitude. A sufficient condition for
retaining minimum normalized attack energy $\lambda_{\min}$ in direction $u$
is
\begin{equation}
    \epsilon^2 u^T
    \left(S_k^0+\Delta S_{\max}\right)^{-1}u
    \geq \lambda_{\min}.
\end{equation}
This condition makes explicit the tradeoff between benign covariance
adaptation and security sensitivity: faster or less bounded adaptation can
improve nominal calibration while expanding the set of attacks whose
normalized evidence falls below threshold.

\begin{figure}[t]
\centering
\fbox{\begin{minipage}{0.92\linewidth}
\centering
\scriptsize
Attacked GPS offset $a_k$\\
$\downarrow$\\
larger GPS residual / innovation $\nu_k$\\
$\downarrow$\\
residual-coupled covariance proposal $\widehat{Q}_k$ increases\\
$\downarrow$\\
$P_k^-$ and $S_k$ increase\\
$\downarrow$\\
normalized attack energy $a_k^TS_k^{-1}a_k$ decreases\\
$\downarrow$\\
alarm can be delayed or avoided
\end{minipage}}
\caption{Covariance-poisoning feedback mechanism. GPS-independent and evidence-gated variants are designed to interrupt or bound this feedback path.}
\label{fig:covariance_causal}
\end{figure}

\section{Baselines}
\label{sec:baselines}

Each detector is calibrated using benign complete runs only. The final
threshold is frozen from all 20 design-corpus runs, and leave-one-run-out
false alarms are reported as a design-corpus diagnostic. Attack data are not
used for threshold selection.

\begin{table*}[t]
\centering
\caption{Detector and estimator variants. Each comparator is calibrated at the same complete-run false-alarm procedure.}
\label{tab:baselines}
\scriptsize
\begin{tabular}{L{0.08\linewidth} L{0.18\linewidth} L{0.24\linewidth} L{0.17\linewidth} L{0.22\linewidth}}
\toprule
\textbf{ID} & \textbf{Score} & \textbf{State update and adaptation} & \textbf{Alarm rule} & \textbf{Attack influence path} \\
\midrule
B1 & Consecutive GPS displacement & No EKF adaptation required for score & One-of-one threshold & Directly sees abrupt jumps, weak for slow drift \\
B2 & Raw GPS-to-prediction residual & Fixed prediction covariance & Three-of-five threshold & GPS affects score but not covariance proposal \\
B3 & NIS & Fixed $Q$ and $R$ EKF & Three-of-five threshold & GPS affects innovation only \\
B4 & NIS & Fixed EKF; reject update above robust innovation gate & Three-of-five threshold & Large GPS innovations rejected from state update \\
B5 & NIS & Fixed EKF; Huber downweights large innovations & Three-of-five threshold & GPS update influence is bounded \\
B6 & Whitened-innovation CUSUM & Fixed EKF; sequential accumulation & One-of-one CUSUM threshold & Persistent residual bias accumulates over time \\
B7 & NIS & Adaptive $Q$ includes GPS residual feedback & Three-of-five threshold & GPS can influence covariance proposal \\
B8 & NIS & Adaptive $Q$ excludes GPS coordinate residual & Three-of-five threshold & GPS affects score but not process features \\
B9 & NIS & Residual feedback accepted only through protected evidence gate & Three-of-five threshold & GPS-only attacks cannot open the gate directly \\
B10 & NIS & Innovation-matching adaptive EKF comparator & Three-of-five threshold & Innovation statistics can influence covariance scale \\
Oracle & NIS & Uses clean-run frozen covariance schedule & Three-of-five threshold & Counterfactual reference, not deployable \\
\bottomrule
\end{tabular}
\end{table*}

\section{Experimental Method}
\label{sec:method}

\subsection{Benign Corpus}

The accepted design corpus contains 20 physical UGV01 runs:
\begin{equation}
2\ \text{speeds} \times 2\ \text{surfaces} \times 5\ \text{trials}
= 20\ \text{runs}.
\end{equation}
The surfaces are smooth kitchen floor and rough permeable concrete. The route
is a 0.5 m square repeated three times. All accepted physical runs use baseline
Wi-Fi. Buffered delay and jitter are applied at replay time without changing
source telemetry or rover behavior.

\subsection{Attack Profiles}

The attack campaign includes step bias, coordinate freeze, value replay, slow
drift, and strategic drift. Step attacks use magnitudes
\[
\{0.5,1,2,3,5,7.5,10\}\ \mathrm{m}
\]
in along-track and cross-track directions. Slow drift uses rates
\[
\{0.01,0.03,0.05\}\ \mathrm{m/s}
\]
in both directions. Attacks start at 25 percent, 50 percent, and 70 percent of
the post-motion run horizon.

\begin{table}[t]
\centering
\caption{Distribution of frozen attack profiles and detector-run evaluations.}
\label{tab:attack_distribution}
\scriptsize
\begin{tabular}{L{0.32\linewidth} r r}
\toprule
\textbf{Attack type} & \textbf{Profiles} & \textbf{Detector-run evaluations} \\
\midrule
Step bias: 7 magnitudes $\times$ 2 directions & 14 & 9,240 \\
Slow drift: 3 rates $\times$ 2 directions & 6 & 3,960 \\
Strategic drift: 1 rate $\times$ 2 directions & 2 & 1,320 \\
Coordinate freeze & 1 & 660 \\
Value replay & 1 & 660 \\
\midrule
Total & 24 & 15,840 \\
\bottomrule
\end{tabular}
\end{table}

\begin{table}[t]
\centering
\caption{Expanded replay grid used for sensitivity analysis.}
\label{tab:expanded_grid}
\scriptsize
\begin{tabular}{lr}
\toprule
\textbf{Quantity} & \textbf{Value} \\
\midrule
Physical benign runs & 20 \\
Detector variants & 11 \\
Attack profiles & 44 \\
Attack start fractions & 3 \\
Unique attack-run-start combinations & 2,640 \\
Detector-run evaluations & 29,040 \\
Invalid condition groups & 0 \\
\bottomrule
\end{tabular}
\end{table}

\subsection{Metrics}

Each replay produces both per-update detector fields and run-level summaries.
For benign runs, the principal metric is complete-run false alarm. For attacks,
the study reports run detection probability, median detection delay, maximum
paired state deviation before alarm, time above the mission tolerance before
alarm, and tolerance-exceeding paired divergence before alarm. A replay
evaluation is counted as tolerance-exceeding if the paired attacked replay
exceeds the selected tolerance before an alarm is raised. The current primary
tolerance is $\epsilon_{\mathrm{req}}=5$ m, with additional post-campaign
tables at 0.25, 0.5, 1, and 2 m.

The paired clean replay is used as the reference for offline state impact:
\begin{equation}
    e_k^{\mathrm{paired}} =
    \|\hat{x}_{k,\mathrm{attack}}-\hat{x}_{k,\mathrm{clean}}\|_2.
\end{equation}
This is not independent physical ground truth; it isolates the effect of the
injected GPS sequence on the same physical log. Confidence intervals are
computed with stratified physical-run-cluster bootstrap resampling so that the
three attack starts within the same physical run are not treated as independent
physical experiments.

\begin{table}[t]
\centering
\caption{Completed offline campaign size with corrected terminology.}
\label{tab:campaign_size}
\scriptsize
\begin{tabular}{lr}
\toprule
\textbf{Quantity} & \textbf{Value} \\
\midrule
Physical benign runs & 20 \\
Detector variants & 11 \\
Attack profiles & 24 \\
Attack start fractions & 3 \\
Unique attack-run-start combinations & 1,440 \\
Expected detector-run evaluations & 15,840 \\
Observed detector-run evaluations & 15,840 \\
Invalid condition groups & 0 \\
\bottomrule
\end{tabular}
\end{table}

Table~\ref{tab:campaign_size} is the frozen, preregistered-style campaign used
for the main reported curves. Table~\ref{tab:expanded_grid} is a replay-only
sensitivity extension with finer step magnitudes and additional drift rates.
The expanded grid is useful for estimating transition behavior, but it is kept
separate from the frozen campaign so that exploratory sensitivity analysis does
not silently replace the main operating point.

\section{Results}
\label{sec:results}

The expanded campaign artifacts were generated under the
\texttt{results/} folder with the same structure as the repository analysis
output directory. The campaign validation file reports a pass: all 1,440
unique attack-run-start combinations and 15,840 detector-run evaluations were
observed, with no invalid condition groups. The 15,840 evaluations are
clustered replay evaluations, not independent physical attack trials.
Unless explicitly labeled as expanded-grid sensitivity analysis, the figures
in this section use the frozen 24-profile campaign. The expanded-grid figures
are stored under \texttt{results/expanded\_grid/} and are used to refine the
transition-region interpretation without replacing the frozen operating point.

\subsection{Result Summary}

The offline campaign supports four main observations. First, the replay and
baseline machinery is complete: all expected detector-run evaluations execute with matched
false-alarm calibration. Second, abrupt step attacks are detectable only when
the offset is large relative to the route scale and GPS uncertainty. Third,
slow drift is substantially harder than abrupt offset: the representative
0.05 m/s cross-track drift is not detected by any variant at the locked
operating point. Fourth, residual-coupled adaptive covariance shows a larger
tolerance-exceeding paired-divergence probability under slow drift than several fixed and
robust baselines, supporting the covariance-poisoning mechanism.

\subsection{False-Alarm Calibration}

Table~\ref{tab:thresholds} summarizes the benign threshold lock. Every
comparator produced one leave-one-run-out alarm out of 20 benign design-corpus
runs. The point estimate is therefore
\begin{equation}
    \hat{P}_{FA} = 1/20 = 0.05,
\end{equation}
with Wilson 95 percent interval approximately $[0.009,0.236]$. This is useful
as a matched design-corpus diagnostic, but it is not yet prospective
false-alarm validation.

\begin{table}[t]
\centering
\caption{Benign threshold-lock summary.}
\label{tab:thresholds}
\scriptsize
\begin{tabular}{L{0.45\linewidth} r r}
\toprule
\textbf{Variant} & \textbf{Threshold} & \textbf{LORO FA} \\
\midrule
Fixed NIS & 16.9015 & 1/20 \\
Naive adaptive & 28.0716 & 1/20 \\
Frozen clean oracle & 28.0716 & 1/20 \\
GPS-independent adaptive & 29.2839 & 1/20 \\
Evidence-gated adaptive & 28.7934 & 1/20 \\
GPS jump & 5.0264 m & 1/20 \\
Raw DT residual & 7.3078 m & 1/20 \\
Robust innovation gate & 25.3239 & 1/20 \\
Huber EKF & 6.4993 & 1/20 \\
CUSUM whitened innovation & 190.3806 & 1/20 \\
Innovation-matching adaptive & 10.2770 & 1/20 \\
\bottomrule
\end{tabular}
\end{table}

The wide confidence interval is important. The point estimate satisfies the
pre-registered target $P_{FA}\leq0.05$, but only 20 physical benign runs are
available. Therefore, the evidence is best described as matched calibration on
the design corpus, not proof that the deployed system will satisfy the same
false-alarm rate on future unseen routes.

\subsection{Step-Attack Detectability}

Figure~\ref{fig:step_detection} shows run-level detection probability for
along-track and cross-track step attacks. The strongest abrupt-offset baseline
is the GPS-jump detector, which reaches approximately
$\epsilon_{90}=7.05$--$7.08$ m. Fixed NIS and raw-residual methods require
approximately 9.7--10 m for high detection probability. Several adaptive
variants do not reach 90 percent detection within the tested 10 m grid.

This result answers RQ1 with a capability limit rather than a strong
protection claim. The detector family can identify large coordinate steps, but
the smallest reliable step is much larger than the 0.5 m square route. This
means that under the current GPS receiver, route, and false-alarm setting, the
system has low sensitivity to small semantic coordinate shifts. The negative
result is still useful because it quantifies the detection boundary rather
than assuming that an EKF innovation test is sufficient.

\begin{figure}[t]
\centering
\includegraphics[width=\linewidth]{results/step_detection_probability.png}
\caption{Frozen-campaign step-attack detection probability for along-track and cross-track GPS offsets.}
\label{fig:step_detection}
\end{figure}

\begin{figure}[t]
\centering
\includegraphics[width=\linewidth]{results/expanded_grid/step_detection_probability.png}
\caption{Expanded-grid step-attack sensitivity analysis with additional replay-only magnitudes around the observed transition region.}
\label{fig:step_detection_expanded}
\end{figure}

\subsection{Slow-Drift Outcomes}

Figure~\ref{fig:drift_outcomes} summarizes slow-drift results. The most
important result is negative but scientifically useful: for the representative
0.05 m/s cross-track drift condition, run-level detection probability is zero
for all variants. However, tolerance-exceeding paired-divergence probability differs across
variants. Naive adaptive covariance reaches approximately 0.150 at the 5 m tolerance, while several
fixed or robust baselines are approximately 0.017.

This result answers RQ2 and RQ3. Slow drift is not merely a smaller version of
the step attack; it interacts with the persistence alarm and covariance
adaptation over time. The detector sees each individual innovation as
plausible enough that no run-level alarm occurs, but the paired state can
still diverge before the alarm. This is the principal security concern for
adaptive uncertainty: the estimator can become less sensitive while still
appearing statistically calibrated on benign data.

\begin{table}[t]
\centering
\caption{Representative 0.05 m/s cross-track drift result at 5 m paired-divergence tolerance.}
\label{tab:drift_5m_tolerance}
\scriptsize
\begin{tabular}{L{0.46\linewidth} c}
\toprule
\textbf{Variant} & \textbf{Tolerance-exceeding probability} \\
\midrule
B7 naive adaptive & 0.150 [0.100, 0.217] \\
Frozen clean oracle & 0.133 [0.083, 0.183] \\
B8 GPS-independent adaptive & 0.100 [0.083, 0.133] \\
B9 evidence-gated adaptive & 0.100 [0.083, 0.133] \\
B3 fixed NIS & 0.017 [0.000, 0.050] \\
B1 GPS jump & 0.017 [0.000, 0.050] \\
B2 raw DT residual & 0.017 [0.000, 0.050] \\
B4 robust innovation gate & 0.017 [0.000, 0.050] \\
B5 Huber EKF & 0.017 [0.000, 0.050] \\
B6 CUSUM whitened innovation & 0.017 [0.000, 0.050] \\
Innovation-matching adaptive & 0.017 [0.000, 0.050] \\
\bottomrule
\end{tabular}
\end{table}

The paired detector-difference table gives the same result as a direct effect
size. For 0.05 m/s cross-track drift, naive adaptive exceeds fixed NIS by
0.133 [0.083, 0.183] in 5 m paired-divergence probability, exceeds robust
innovation gating by 0.133 [0.083, 0.183], and exceeds Huber EKF by 0.133
[0.100, 0.183]. Compared with evidence-gated adaptation, the difference is
smaller but still positive: 0.050 [0.017, 0.100]. These paired differences are
computed with physical-run-clustered bootstrap resampling.

\begin{figure}[t]
\centering
\includegraphics[width=\linewidth]{results/drift_attack_outcomes.png}
\caption{Frozen-campaign detection probability and tolerance-exceeding paired-divergence probability under slow drift.}
\label{fig:drift_outcomes}
\end{figure}

\begin{figure}[t]
\centering
\includegraphics[width=\linewidth]{results/expanded_grid/drift_attack_outcomes.png}
\caption{Expanded-grid drift sensitivity analysis with additional replay-only drift rates.}
\label{fig:drift_outcomes_expanded}
\end{figure}

\subsection{Covariance-Poisoning Evidence}

Figure~\ref{fig:cov_poisoning} compares covariance response and paired state
impact under drift. The naive residual-coupled adaptive method can increase
the covariance trace under attack, which reduces normalized innovation
sensitivity and increases the amount of paired state deviation that can occur
before alarm. This supports the mechanism in
Theorem~\ref{thm:dynamic_poisoning}: adaptive uncertainty can become a
security-relevant feedback path when attacker-controlled GPS residuals
influence covariance proposals.

The observed covariance response should be interpreted as evidence for the
mechanism, not as final proof of an operational attack advantage in every
condition. The strongest defensible statement is that attacked
GPS-residual-dependent adaptation can inflate uncertainty and suppress
normalized evidence. A stronger claim, that covariance poisoning reliably
increases mission-relevant physical harm relative to all clean-covariance
controls, requires independent physical ground truth and prospective
validation.

\begin{table}[t]
\centering
\caption{Gate and covariance behavior for 0.05 m/s cross-track drift.}
\label{tab:gate_behavior}
\scriptsize
\begin{tabular}{L{0.30\linewidth} c c c}
\toprule
\textbf{Variant} & \textbf{Residual feedback} & \textbf{Ind. evidence} & \textbf{Attack-window $Q$ ratio} \\
\midrule
Fixed NIS & 0.000 & 0.366 & 1.000 \\
Naive adaptive & 1.000 & 0.366 & 1.040 \\
GPS-independent adaptive & 0.000 & 0.366 & 1.000 \\
Evidence-gated adaptive & 0.366 & 0.366 & 1.025 \\
\bottomrule
\end{tabular}
\end{table}

Table~\ref{tab:gate_behavior} explains the mechanism behind the variants.
Naive adaptive keeps residual feedback active throughout the attack window,
whereas GPS-independent adaptation removes that path. Evidence gating permits
feedback only when protected IMU or timing evidence is present, which occurs
for about 36.6 percent of the representative drift windows. The corresponding
attack-window $Q$ trace ratio is 1.040 for naive adaptive and 1.025 for
evidence-gated adaptation, indicating that the gate reduces but does not
eliminate adaptive covariance growth.

\begin{figure}[t]
\centering
\includegraphics[width=\linewidth]{results/covariance_poisoning.png}
\caption{Frozen-campaign covariance and paired-state response under slow drift.}
\label{fig:cov_poisoning}
\end{figure}

\subsection{Buffered Transport Check}

A deterministic edge-side buffered condition with 200 ms delay and 40 ms
jitter was applied to the diagnostic trial-5 subset. Source packets and rover
behavior were unchanged. All 11 variants produced zero benign buffered alarms
on this diagnostic subset.

The buffered result is a sanity check rather than a complete timing study. It
shows that the frozen detectors do not immediately false alarm under the
diagnostic replayed delay/jitter condition. A final edge-performance claim
will require running the frozen detector live and reporting source-to-estimate
and source-to-alarm latency, update rate, queue depth, and missed-deadline
fraction.

\subsection{Post-Campaign Sensitivity, Provenance, and Runtime}

The repository includes a post-campaign analysis pass that does not replay logs
or require new data. It derives multiple paired-divergence thresholds,
paired detector differences with physical-run-clustered bootstrap intervals,
gate behavior summaries, and provenance hashes from the existing campaign CSVs.
These outputs are intended to supplement the main figures with numeric tables:
\begin{itemize}
    \item tolerance-exceeding paired divergence before alarm at 0.25, 0.5,
    1, 2, and 5 m;
    \item paired detector differences such as naive adaptive minus fixed NIS,
    and naive adaptive minus evidence-gated adaptation;
    \item mean residual-feedback activation, independent-evidence fraction,
    and covariance trace ratios by attack and detector;
    \item hashes for the manifest, campaign summary, threshold policy, and
    uncertainty-policy configuration.
\end{itemize}

The post-campaign provenance file hashes nine principal artifacts, including
the benign manifest, campaign summary, aggregate results, epsilon summary,
validation certificate, locked thresholds, and attack and uncertainty policy
files. This makes the reported tables traceable to a specific analysis state.
The runtime benchmark uses existing logs and reports steady-state replay
runtime of approximately 2.3--2.6 ms per update for the detector variants after
warm-up. The first measured fixed-covariance replay includes one warm-up
outlier and is not used as the steady-state runtime claim. These timings are
offline Python replay estimates; a final live edge deployment must still
measure CPU load, memory use, source-to-estimate latency, and source-to-alarm
latency during streaming operation.

\section{Discussion}
\label{sec:discussion}

The present results support a careful claim: adaptive uncertainty can be a
security-relevant attack surface in an edge-hosted mobile digital twin. The
evidence is strongest for the mechanism. The paired replay shows covariance
inflation and reduced normalized sensitivity under slow GPS drift, and the
naive adaptive variant has higher tolerance-exceeding paired-divergence probability than
several fixed and robust baselines. The paired differences are largest against
fixed, robust-gated, Huber, and CUSUM baselines, and smaller against
GPS-independent and evidence-gated adaptation. This pattern is consistent with
the proposed feedback story: removing or gating GPS-residual feedback reduces
the direct attack-to-covariance path, but does not by itself make slow drift
easy to detect.

The results do not yet support the stronger claim that the proposed
evidence-gated defense solves GPS manipulation. Slow drift remains difficult
for every tested detector under the current false-alarm requirement and
platform uncertainty. This should be framed as a security limit, not hidden as
a failed positive result.

The step-attack results also show an important capability limit. Detection
becomes reliable only for large offsets, especially relative to the 0.5 m
square route. This suggests that for low-cost GPS and indoor or partially
obstructed operation, small semantic coordinate attacks may be statistically
hard to separate from benign uncertainty without additional independent
reference information.

\begin{table*}[t]
\centering
\caption{Claims supported by the current evidence and remaining limitations.}
\label{tab:claims_evidence}
\scriptsize
\begin{tabular}{L{0.22\linewidth} L{0.37\linewidth} L{0.32\linewidth}}
\toprule
\textbf{Claim} & \textbf{Current evidence} & \textbf{Limitation} \\
\midrule
Real-log security replay is reproducible &
20 physical runs, 11 variants, 24 frozen attack profiles, three starts,
15,840/15,840 frozen detector-run evaluations, and 29,040/29,040 expanded-grid
evaluations completed &
Offline replay only; not yet live attack execution \\
Matched false-alarm comparison is implemented &
Every variant calibrated with benign complete runs; LORO result is 1/20 for
all variants &
20-run design corpus gives a wide confidence interval \\
Large step attacks are detectable &
GPS-jump reaches $\epsilon_{90}\approx7.05$--$7.08$ m; fixed/raw residual
approach high detection near 10 m &
Small 0.5--2 m manipulations remain weakly detected \\
Slow drift is difficult &
0.05 m/s cross-track drift has zero run-level detection probability across
variants &
Result is counterfactual replay against logs, not live physical spoofing \\
Covariance poisoning mechanism exists &
Naive adaptive covariance shows increased tolerance-exceeding paired-divergence probability
and covariance response under drift; paired difference versus fixed is 0.133
[0.083, 0.183] for the representative 0.05 m/s cross-track drift &
Operational physical harm needs independent trajectory ground truth \\
Evidence gating is promising but not solved &
Gating reduces residual-feedback activation from 1.000 to 0.366 and lowers the
representative attack-window $Q$ ratio from 1.040 to 1.025 &
Current results do not show decisive dominance over all baselines \\
Reproducibility package exists &
Manifest, campaign, threshold, and policy artifacts are hashed; single-command
campaign, post-campaign, threshold-sweep, and runtime scripts are documented &
External replication still requires packaging raw logs and environment versions \\
\bottomrule
\end{tabular}
\end{table*}

\section{Limitations and Future Work}
\label{sec:limitations}

The current campaign is an offline counterfactual replay study. Attacks are
injected into real logs, not executed live against the rover. State impact is
measured against paired clean replay rather than independent physical ground
truth. The 20 benign runs were available during alarm design, so
leave-one-run-out false-alarm results are design-corpus diagnostics rather
than prospective validation. A final strong empirical claim requires
independent camera or AprilTag ground truth, a fresh prospective corpus, and
live edge execution with frozen policies.

The expanded attack grid and threshold sweep improve resolution around the
observed transition regions, but they remain replay-only sensitivity analyses
on the same 20 physical runs. They should not be counted as new physical
experiments. Similarly, the post-campaign paired-difference tables strengthen
statistical interpretation, but they do not replace prospective validation.

Future work will add fixed-camera AprilTag or ChArUco tracking, collect an
untouched prospective validation set, and compare physical trajectory error,
GPS error, encoder/dead-reckoning error, and digital-twin reported error under
the frozen detector.

\section{Reproducibility Package}
\label{sec:reproducibility}

The artifact package contains raw benign telemetry logs, a benign manifest
with SHA-256 hashes, frozen threshold and uncertainty-policy files, analysis
scripts, generated CSV summaries, plots, and a data dictionary. The full
offline baseline campaign is regenerated with:
\begin{verbatim}
python -m DigitalTwin.analysis.real_data_study \
  --bootstrap-iterations 2000
\end{verbatim}
The post-campaign tables are regenerated without replaying logs with:
\begin{verbatim}
python -m DigitalTwin.analysis.post_campaign_analysis \
  --bootstrap-iterations 2000
\end{verbatim}
ROC-style threshold-sweep curves require replaying scores and are therefore
separate:
\begin{verbatim}
python -m DigitalTwin.analysis.threshold_sweep \
  --expanded-attacks --grid-points 80
\end{verbatim}
Combined content-and-buffered-delay replay is generated with:
\begin{verbatim}
python -m DigitalTwin.analysis.real_data_study \
  --include-buffered-attacks \
  --bootstrap-iterations 2000
\end{verbatim}
The expanded replay grid is regenerated with:
\begin{verbatim}
python -m DigitalTwin.analysis.real_data_study \
  --expanded-attack-grid \
  --bootstrap-iterations 2000
\end{verbatim}
Offline runtime cost is estimated with:
\begin{verbatim}
python -m DigitalTwin.analysis.runtime_benchmark \
  --max-runs 4
\end{verbatim}
The expanded-grid, threshold-sweep, and combined buffered-attack commands can be long-running
because they replay many detector and attack combinations.

\section{Conclusion}

This paper presents a real-log security study of adaptive uncertainty in an
edge-hosted digital twin for a mobile rover. The system integrates UGV01
firmware telemetry, BN220 GPS, edge timing, EKF replay, semantic GPS attack
injection, matched false-alarm thresholding, and 11 detector baselines. The
completed campaign validates the offline infrastructure over 15,840 detector-run
evaluations, with an expanded 29,040-evaluation replay grid and post-campaign
paired comparisons. The results show that large GPS jumps are detectable but
slow drift is hard, and that naive residual-coupled covariance adaptation can
increase tolerance-exceeding paired divergence before alarm. The strongest
defensible conclusion is therefore not that the detector solves GPS spoofing.
Instead, the paper shows that uncertainty adaptation changes the attack
surface: it can improve operational calibration while simultaneously reducing
normalized evidence for patient semantic manipulations. This motivates future
edge digital twins that separate security prediction from attacked
measurements, constrain adaptive feedback with independent evidence, and
report detectability limits rather than only binary attack decisions.

\appendices
\section{Implementation Notes}

The active firmware tree used for the UGV01 experiments is
\texttt{ugv01\_gps\_dev}. The BN220 receiver remains on the verified original
connection: white signal wire to UGV01 RX, red wire to 5 V, and black wire to
ground. Command \texttt{T:146} is the GPS-only telemetry request and
\texttt{T:147} is the combined telemetry request used for the real-log replay
study. These details are included for reproducibility and are not assumptions
in the attack-detection theory.

\section*{Acknowledgment}

Funding and advisor acknowledgments can be added here.

\begin{thebibliography}{99}

\bibitem{ref2}
W. Aldosari, ``Deep learning-based location spoofing attack detection and
time-of-arrival estimation through power received in IoT networks,''
\emph{Sensors}, vol. 23, no. 23, 2023.

\bibitem{ref7}
S. Cai, W. Gou, W. Wen, X. Lu, J. Fan, and X. Guo, ``Design and development
of a low-cost UGV 3D phenotyping platform with integrated LiDAR and electric
slide rail,'' \emph{Plants}, vol. 12, no. 3, 2023.

\bibitem{ref11}
R. Gonzalez and K. Iagnemma, ``Slippage estimation and compensation for
planetary exploration rovers: State of the art and future challenges,''
\emph{Journal of Field Robotics}, vol. 35, no. 4, pp. 564--577, 2017.

\bibitem{ref43}
G. Pettorru, V. Pilloni, and M. Martalo, ``Trustworthy localization in IoT
networks: A survey of localization techniques, threats, and mitigation,''
\emph{Sensors}, vol. 24, no. 7, 2024.

\bibitem{ref48}
M. J. C. S. Reis and A. J. D. Reis, ``Edge-based real-time fault detection in
UAV systems via B-spline telemetry reconstruction and lightweight hybrid AI,''
\emph{Sensors}, vol. 25, no. 16, 2025.

\bibitem{ref60}
K. S. Tharayil \emph{et al.}, ``Sensor defense in-software: Practical
software based detection of spoofing attacks on position sensor,'' arXiv
preprint arXiv:1905.04691, 2019.

\bibitem{ref61}
C. Toma, M. Popa, B. Iancu, M. Doinea, A. Pascu, and F. Ioan-Dutescu, ``Edge
machine learning for automated decision and visual computing of robots, IoT
embedded devices or UAV-drones,'' \emph{Electronics}, vol. 11, no. 21, 2022.

\bibitem{ref65}
Z. Yu, Z. Kaplan, Q. Yan, and N. Zhang, ``Security and privacy in the emerging
cyber-physical world: A survey,'' \emph{IEEE Communications Surveys \&
Tutorials}, vol. 23, no. 3, pp. 1879--1919, 2021.

\bibitem{mo2010}
Y. Mo and B. Sinopoli, ``False data injection attacks in control systems,'' in
\emph{Proc. 1st Workshop on Secure Control Systems}, 2010.

\bibitem{pasqualetti2013}
F. Pasqualetti, F. Dorfler, and F. Bullo, ``Attack detection and
identification in cyber-physical systems,'' \emph{IEEE Transactions on
Automatic Control}, vol. 58, no. 11, pp. 2715--2729, 2013.

\bibitem{fawzi2014}
H. Fawzi, P. Tabuada, and S. Diggavi, ``Secure estimation and control for
cyber-physical systems under adversarial attacks,'' \emph{IEEE Transactions
on Automatic Control}, vol. 59, no. 6, pp. 1454--1467, 2014.

\bibitem{page1954}
E. S. Page, ``Continuous inspection schemes,'' \emph{Biometrika}, vol. 41,
no. 1/2, pp. 100--115, 1954.

\bibitem{huber1981}
P. J. Huber, \emph{Robust Statistics}. Wiley, 1981.

\end{thebibliography}

\end{document}
